Integrální manometrický permeametr, jehož konstrukcí a~parametry se detailně zabývá Kapitola \ref{s:integralni_manometricky_permeametr}, lze pro účely matematického modelování zredukovat na dva objekty: membránu a~permeační celu. Děje probíhající uvnitř membrány budeme označovat jako \emph{vnitřní}: jde o~děje na mikroskopické až molekulární úrovni, kde zásadní budou úvahy vycházející z termodynamických principů. Děje probíhající v~permeační cele budeme označovat jako \emph{vnější}: jde o~děje makroskopické, jejichž projevy jsou měřitelné běžně dostupnou technikou, jako jsou teplota či tlak.

V Kapitole \ref{s:vnitrni_dynamika_systemu} se věnujeme vnitřní dynamice: odvozujeme řídicí rovnici vybraných stavových veličin jako funkcí prostoru a~času. Dostáváme nelineární parciální diferenciální rovnici známou jako \emph{Porous Medium Equation} (dále PME), patřící do širší třídy rovnic nelineární difuse (\emph{Nonlinear Diffusion Equation}). Pokračujeme nalezením jejího asymptotického řešení pro Dirichletovu úlohu, která je přirozeným popisem experimentálního uspořádání. Tok plynu membránou se ukáže být prostorovou invariantou, jak lze předjímat i ze zákonů zachování. Postupnou eliminací nejprve časové dimenze použitím asymptotického řešení, následovaném eliminací prostorové dimenze nalezením prostorové invarianty, získáváme vyjádření toku plynu membránou jako funkci okrajových podmínek. 

V Kapitole \ref{s:vnejsi_dynamika_systemu} uvažujeme látkovou bilanci permeační cely, kde vstupem je, na základě principu lokální rovnováhy, dříve zjištěný tok plynu membránou. Získáváme časový průběh tlaku v~permeační cele popsaný obyčejnou nelineární diferenciální rovnicí prvního řádu; počáteční podmínka odpovídá stavu permeační cely před zahájením experimentu. Tlak v~zásobní cele lze předepsat buď konstantou, nebo daty poskytnutými regulátorem; v~případě konstanty lze zcela analogicky uvažovat o~odezvě systému na Heavisideovu funkci, resp. její násobek. Analytické řešení počáteční úlohy lze nalézt pouze implicitně, neboť se neobejde bez využití nekonečných řad v~podobě hypergeometrických funkcí; pro vybrané hodnoty parametru však dokážeme najít i vyjádření explicitní.

V Kapitole \ref{s:identification} vyvíjíme metodiku pro stanovení hodnot dvou neznámých fyzikálních parametrů, jednoho multiplikativního a~jednoho exponenciálního charakteru, kde studujeme zejména druhý z nich. Nejprve vyšetřujeme průběh přechodové charakteristiky a vliv exponenciálního parametru, rekurentní charakter zjištěných vztahů dodá problematice interpretaci s použitím statistických momentů. Dále odvozujeme dvě nezávislé metody experimentální identifikace. První metoda stanovuje exponenciální parametr na základě křivosti odezvy systému na Heavisideovu funkci, kde využívá konceptu časové konstanty systému prvního řádu jako normalizačního mechanismu, hodnota parametru pak je jednoznačně určena hodnotou odezvy systému v~konkrétním normalizovaném čase. Druhá metoda stanovuje hodnotu parametru na základě porušení principu superpozice, konkrétně nezávislosti na měřítku (\emph{scale-invariance}): exponenciální parametr určujeme kvantifikací odchýlení od přímé úměry závislosti počáteční směrnice odezvy systému na amplitudě vstupního signálu.